% ============================================================
% SECTION 3: EVALUATION PROTOCOL. Written clean 2026-07-26.
% Title chosen by Mohammad 2026-07-26 over "What Honest Evaluation Requires":
%   "honest" implies the alternative is dishonesty, which is not this paper's
%   position. Satuluri et al. stated the requirement correctly in 2011.
% Weight: MEDIUM-HEAVY (PAPER_STRUCTURE.md). This is the instrument.
%
% Evidence, control by control:
%   transfer scoring   TERMINOLOGY.md; exp_P SUMMARY V2 (45/63, 60/63, 15/63);
%                      exp_G SUMMARY Table A (uniform, 6/6 sign flips)
%   granularity        exp_P finding 4 (com-DBLP Louvain 0.1930 -> 0.2996 at
%                      gamma=576); exp_L (Leiden 0.1930 / 0.2971)
%   chance floor       exp_P finding 4 (com-Amazon 0.0086@k=231 -> 0.0567@k=8651)
%   compute matching   exp_P V1/V3 and caveat C3; exp_C section 1 (2 restarts
%                      in all 48 cells); exp_AA (budget buys 1 restart)
%   configuration null exp_B SUMMARY (17 networks; dQ_fixed and delta tables)
%   permuted weights   exp_V V2 (email-Enron; ca-CondMat rewiring)
%   realized retention exp_C section 0 (sampler audit); exp_AB section 1 (floors)
%   end-to-end cost    exp_P V5 (0.66x max among quality-preserving cells)
%
% NOTE: the modularity definition sits here because the m-dependence of its
% null term IS the argument in 3.1. If Section 2 later defines it, delete
% Eq. (1) here and cross-reference instead.
% STYLE: no em-dashes, no AI-characteristic phrasing (see STYLE.md).
% ============================================================

Section~\ref{sec:intro} listed four requirements, each belonging to a specific
class of method rather than to sparsification in general. This section states
them as controls. Each is given with the failure it catches and with what we
measure when it is absent. Table~\ref{tab:protocol} collects them. Every number
in the sections that follow is produced under all of them.

\subsection{One yardstick for two partitions}\label{sec:protocol:transfer}

Detection may use any graph; scoring must use the original one. Let
$\mathcal{A}$ be a detection algorithm, $G$ the original graph and $G'$ a
sparsified copy of it. The two partitions under comparison are
$P_{\mathrm{sparse}} = \mathcal{A}(G')$ and $P_{\mathrm{base}} = \mathcal{A}(G)$,
and since both are partitions of the same vertex set, the question the pipeline
raises is which of them better describes $G$. Two accountings of that question
appear in this literature:
\begin{align}
\Delta Q_{\mathrm{transfer}} &= Q(P_{\mathrm{sparse}}, G) - Q(P_{\mathrm{base}},
G), \label{eq:transfer}\\
\Delta Q_{\mathrm{self}} &= Q(P_{\mathrm{sparse}}, G') - Q(P_{\mathrm{base}}, G).
\label{eq:self}
\end{align}
They differ in one place, the graph on which the first term is evaluated.
Equation~\eqref{eq:transfer} measures both partitions with a single instrument.
Equation~\eqref{eq:self} changes the instrument between the arm under test and
the arm it is tested against, and the instrument it applies to the former is the
graph from which edges have been removed. We report
$\Delta Q_{\mathrm{transfer}}$ throughout; $\Delta Q_{\mathrm{self}}$ appears in
released benchmark tooling and in recent work in this line, documented in
Section~\ref{sec:related}.

The two accountings diverge systematically, and the objective shows why.
Modularity scores a partition $P$ of a graph with $m$ edges by
\begin{equation}\label{eq:modularity}
Q(P, G) \;=\; \sum_{C \in P} \left[ \frac{m_C}{m} - \gamma
\left( \frac{D_C}{2m} \right)^{\!2} \right],
\end{equation}
with $m_C$ the number of edges inside community $C$, $D_C$ the sum of degrees in
$C$, and $\gamma$ a resolution parameter equal to one in the standard definition.
The subtracted term is normalized by $m$. Deleting edges shrinks it, so
modularity computed on a sparsified graph rises whether or not the partition has
improved, and the rise is larger the more edges are removed.

The effect is not small and it is not confined to sparsifiers that select
intelligently. Table~\ref{tab:uniform} scores uniform random edge deletion, whose
selection rule uses no information at all, both ways on three networks at two
retention levels. The sign of the conclusion is opposite in every one of the six
cells, and the gap between the two accountings roughly doubles when retention
falls from $0.9$ to $0.8$, which is the behaviour \eqref{eq:modularity} predicts.

\begin{table}[t]
\centering
\caption{Uniform random edge deletion under the two accountings, Leiden, mean of
three seeds. $\Delta Q_{\mathrm{self}}$ is positive in every cell and
$\Delta Q_{\mathrm{transfer}}$ is negative in every cell. Uniform deletion uses
no structural information, so the inflation comes from the objective rather than
from any selection rule.}
\label{tab:uniform}
\small
\setlength{\tabcolsep}{4pt}
\begin{tabular}{@{}lccrr@{}}
\toprule
Network & $Q_{\mathrm{base}}$ & $\rho$ & $\Delta Q_{\mathrm{self}}$ &
$\Delta Q_{\mathrm{transfer}}$\\
\midrule
ca-CondMat  & 0.729 & 0.800 & $+0.0111$ & $-0.0164$\\
ca-CondMat  & 0.729 & 0.900 & $+0.0053$ & $-0.0064$\\
email-Enron & 0.611 & 0.800 & $+0.0052$ & $-0.0162$\\
email-Enron & 0.611 & 0.900 & $+0.0011$ & $-0.0086$\\
com-DBLP    & 0.826 & 0.800 & $+0.0078$ & $-0.0247$\\
com-DBLP    & 0.826 & 0.900 & $+0.0040$ & $-0.0098$\\
\bottomrule
\end{tabular}
\end{table}

Across a larger grid the pattern holds under every algorithm we tested.
Table~\ref{tab:accounting} covers 63 configurations spanning three detection
algorithms, seven networks and three sparsifier settings. Modularity read off the
sparsified graph is positive in 60 of them and reaches $+0.537$; the same
partitions scored on the original graph are positive in 15. The sign of the
conclusion differs between the two accountings in 45 of the 63 configurations.
Eleven cells exceed a compute-matched baseline, all of them under label
propagation, and they are not a subset of the 15: two of them sit below the plain
baseline mean while beating a matched control that had drawn too few restarts,
which is the failure Section~\ref{sec:protocol:compute} treats.

\begin{table}[t]
\centering
\caption{The two accountings across 63 configurations, seven networks by three
sparsifier settings under each algorithm. A flip is a cell where
$\Delta Q_{\mathrm{self}} > 0$ and $\Delta Q_{\mathrm{transfer}} < 0$, so that
the two accountings disagree on the sign of the conclusion.}
\label{tab:accounting}
\small
\setlength{\tabcolsep}{4pt}
\begin{tabular}{@{}lcccc@{}}
\toprule
 & Infomap & Louvain & Lab.\ prop. & Total\\
\midrule
Cells                              & 21 & 21 & 21 & 63\\
$\Delta Q_{\mathrm{self}} > 0$     & 20 & 20 & 20 & 60\\
$\Delta Q_{\mathrm{transfer}} > 0$ & 6  & 0  & 9  & 15\\
Flips                              & 14 & 20 & 11 & 45\\
max $\Delta Q_{\mathrm{self}}$ & $+0.312$ & $+0.260$ & $+0.537$ & $+0.537$\\
\bottomrule
\end{tabular}
\end{table}

\subsection{Matching granularity}\label{sec:protocol:granularity}

Sparsification fragments a graph, and an algorithm that chooses its own
granularity responds by returning more communities. Chen et
al.~\cite{chen2024demystifying} report community counts rising by three orders of
magnitude across a pruning sweep. Purity, normalized mutual information and the
average best-match $F_1$ over reference communities that we report all increase
with the number of communities, so a recovery comparison between partitions of
different granularity measures that difference along with whatever else it is
meant to measure. Hamann et al.~\cite{hamann2016structure} warned that normalized
mutual information is unsuitable for comparing partitions of sparsified networks,
because it reports higher similarity for larger numbers of
communities~\cite{vinh2010information}, and pointed to the adjusted Rand index
instead; the control below operationalizes that warning.

The control is a partition of the \emph{original} graph at a matched community
count, obtained by raising the resolution parameter $\gamma$ in
\eqref{eq:modularity} until the counts agree. Table~\ref{tab:granularity} applies
it to com-DBLP under Louvain. Local similarity sparsification at realized
retention $0.503$ raises average best-match $F_1$ from $0.102$ to $0.193$, an
apparent gain of $+0.091$, while raising the number of communities of size at
least three from 220 to 10,947. The original graph partitioned at $\gamma = 640$
reaches 11,159 such communities and scores $0.300$. The finer partition is worth
more than twice what the sparsification appeared to be worth, and it requires no
sparsifier. The same holds for degree-based sparsification on the same network,
where the apparent change is negative and the matched control is positive.
Leiden reproduces the sparsified figure to four decimal places and its own
control, at $\gamma = 576$ and 10,812 communities, reaches $0.297$, so the effect
belongs to the modularity objective's default granularity rather than to one
optimizer's search.

\begin{table}[t]
\centering
\caption{Granularity control on com-DBLP under Louvain, average best-match $F_1$
over reference communities of size at least three. $k_{\geq 3}$ counts
communities of that size; $\gamma$ is the resolution used on the \emph{original}
graph. Each sparsified arm is followed by its control, a partition of the
original graph at the same $k_{\geq 3}$. The floor is a size-matched random
partition at the arm's own $k_{\geq 3}$. Baseline and sparsified rows are means
over three to six seeds; controls are single runs.}
\label{tab:granularity}
\small
\setlength{\tabcolsep}{4pt}
\begin{tabular}{@{}lrrcc@{}}
\toprule
Condition & $k_{\geq 3}$ & $\gamma$ & $F_1$ & floor\\
\midrule
Original graph               & 220    & 1   & $0.102 \pm 0.005$ & 0.019\\
DSpar, $\rho = 0.500$        & 1{,}684  & 1   & $0.098 \pm 0.002$ & 0.049\\
\quad matched control        & 1{,}657  & 24  & $0.125$           & \\
L-Spar, $\rho = 0.503$       & 10{,}947 & 1   & $0.193 \pm 0.000$ & 0.090\\
\quad matched control        & 11{,}159 & 640 & $0.300$           & \\
\bottomrule
\end{tabular}
\end{table}

Where the counts cannot be brought together, we draw no recovery conclusion. The
comparability rule is $|\Delta k|/k < 25\%$, and the cells it excludes are
reported as excluded rather than compared.

\subsection{Subtracting chance}\label{sec:protocol:chance}

Recovery is reported with chance-corrected indices. Where an uncorrected measure
is used because the reference communities are overlapping and numerous, as with
average best-match $F_1$ on the largest networks, we also compute the score of a
random partition with the same community size distribution, and report it beside
the measurement.

That floor moves with the quantity the previous control governs. On com-Amazon
under Louvain, a size-matched random partition scores $0.009$ against the
reference communities at 231 communities and $0.057$ at 8,651, a $6.6\times$ rise
produced by fragmentation alone. Any comparison of best-match $F_1$ across
different community counts reads part of that slope, and subtracting the floor is
what separates the two.

\subsection{An equal compute budget}\label{sec:protocol:compute}

Every detection algorithm we use is stochastic, so a single run is a draw rather
than an answer, and a pipeline that sparsifies and then detects has spent more
time than one detection run. The registered control gives the baseline as many
independent restarts as fit the pipeline's wall clock and takes the best of them.

That control is weaker than it appears, and its weakness favours the treatment.
When the pipeline is fast the budget buys very little: two restarts in all 48
cells of one sweep, and between 1 and 16 in another, with 1 in most cells on the
two largest networks. Label propagation shows what this costs. Under the
compute-matched control it is the strongest of the three algorithms in that
sweep, with a mean change of $+0.009$; against the best of 20 plain restarts of
the same baseline it is the weakest, at $-0.106$. On email-Enron its two apparent
gains of $+0.146$ and $+0.161$ become $-0.068$ and $-0.053$, because the
compute-matched baseline had drawn 1 to 2 restarts and never sampled the
algorithm's good mode, whose modularity is $0.523$ against the matched control's
$0.308$.

We therefore report both: the compute-matched figure, which is the fair
accounting of what the pipeline costs, and the best of a fixed number of
restarts, which is a strictly more generous baseline that can only handicap the
treatment. A result counts as surviving only if it survives both. This is also
the reason the granularity control above is applied to results that have already
passed the compute test, rather than in place of it.

\subsection{Nulls}\label{sec:protocol:nulls}

Two of the quantities this literature reads as evidence of community structure
can be produced without any community structure at all. Each has a null that
detects the substitution.

\paragraph{A degree sequence without communities}
A degree-based sparsifier selects on a function of the degree sequence, so an
effect it produces on a heavy-tailed graph may follow from that sequence alone.
The control is a degree-preserving rewiring~\cite{molloy1995critical}, built by
double edge swaps, which holds every vertex degree exactly and destroys the
community structure. A signal that reappears on the rewiring is not evidence
about communities.

Two quantities fail this test. The first is the change in modularity of a fixed
partition when the graph is sparsified, which is positive on all seventeen
networks we tested; on their rewirings it is positive on all seventeen as well,
and larger than the real value on fifteen of them. The second is the separation
between the scores a degree-based sparsifier assigns to within-community and
between-community edges, which on the rewirings is larger than on the real graphs
in sixteen of seventeen cases, with a median ratio of $1.60$. Real community
structure suppresses that separation rather than producing it. Neither quantity
is used as evidence in this paper. The mechanism they belong to is the subject of
Section~\ref{sec:mechanism}, where the same null is what makes the causal test
interpretable.

\paragraph{Weights without their placement}
A related family of proposals keeps every edge and reweights it by a similarity
score. The corresponding null permutes the computed weights uniformly at random
across the edges, which preserves the weight distribution exactly and destroys
every association between a weight and the structure that produced it. We are not
aware of its use in this literature.

It changes the reading of the one modularity gain we obtained from four ways of
deploying a similarity signal. On email-Enron, weighting each edge by
$1 + J(u,v)$, with $J$ the Jaccard similarity of the endpoint neighbourhoods,
raises modularity measured on the original graph to $0.6153 \pm 0.0034$ against a
baseline of $0.6071$ and a compute-matched best of $0.6051$, a gain of $+0.0102$
at five pooled standard deviations. Permuting those same weights across edges
gives $0.6168 \pm 0.0020$, a gain of $+0.0117$, which is larger. On a
degree-preserving rewiring of ca-CondMat, where the similarity signal is absent
by construction, weighting still produces $+0.0023$ and clears the same
significance test. What the weights contribute is heterogeneity in the
optimizer's search rather than information about communities.

\subsection{Realized retention}\label{sec:protocol:retention}

Retention is measured, never assumed. The sampling scheme in common use clips
per-edge probabilities at one, and between 13\% and 35\% of edges reach that
clip, so the expected number of retained edges falls below the nominal $\alpha m$
and saturates: at $\alpha = 1.0$ one of our networks retains $0.662$ of its
edges, and at $\alpha = 0.95$ another retains $0.447$. Three samplers that all
describe themselves by the same $\alpha$ realize $0.33$ to $0.55$, $0.37$ to
$0.68$, and, once the scale factor is solved by bisection, $0.70$ to $0.95$.
Sparsifiers are compared at matched realized retention throughout, and the
realized value appears in every table.

Two properties of the local selection rules limit what can be matched. The
achievable retentions form a coarse grid, so that on com-Amazon local similarity
sparsification can produce only $0.301$ and $0.542$ and nothing between them. And
several methods cannot reach aggressive retention on a sparse graph at all,
because each vertex keeps at least one incident edge; the resulting floors reach
$0.36$ on the sparsest network we test, which is reported with the results rather
than treated as a free parameter.

A comparison must also be like for like in what is removed. A method that deletes
vertices as well as edges scores its partition on a smaller and denser
subpopulation, so the fraction of vertices retained belongs in the table beside
the fraction of edges. Every sparsifier in this study removes edges only, and we
report both fractions.

\subsection{Cost measured end to end}\label{sec:protocol:cost}

The sparsifier is a computation and is charged to the pipeline that uses it, as
it was in the original report~\cite{satuluri2011local} and as Gottesb\"{u}ren et
al.~\cite{gottesburen2025linear} recommend after measuring the same overhead in
graph partitioning. The comparison is the whole pipeline, sparsifier included,
against one run of the same detection algorithm on the original graph.

Among the cells of our largest sweep that preserve quality, the greatest
end-to-end speedup is $0.66\times$, which is slower than not sparsifying.
Speedups above $1.5\times$ occur in 12 of 63 cells, and every one of them loses
modularity. Discounting the sparsifier entirely and timing only the detection
step, the median speedup is between $1.15\times$ and $1.63\times$ by algorithm,
well short of the factor of two that halving the edges suggests, because all
three algorithms return more communities on a sparsified graph and do more work
per pass.

\begin{table*}[t]
\centering
\caption{The protocol. Each control is motivated by a specific failure, and the
last column reports what we measure when the control is absent.}
\label{tab:protocol}
\begin{tabular}{@{}p{0.19\textwidth}p{0.29\textwidth}p{0.42\textwidth}@{}}
\toprule
Failure & Control & Measured consequence of omitting it\\
\midrule
Quality read off the sparsified graph &
Score both partitions on the original graph, Eq.~\eqref{eq:transfer} &
The sign of the conclusion differs in 45 of 63 configurations; positive in 60 of
63 one way and in 15 of 63 the other (\S\ref{sec:protocol:transfer})\\
Finer partitions score higher &
Partition the original graph at a matched community count &
com-DBLP: an apparent $+0.091$ in best-match $F_1$ against a control that reaches
$+0.198$ with no sparsifier (\S\ref{sec:protocol:granularity})\\
Uncorrected agreement measures &
Chance-corrected indices, and a size-matched random-partition floor &
The floor rises $6.6\times$ between 231 and 8,651 communities on one graph
(\S\ref{sec:protocol:chance})\\
Unequal compute &
Compute-matched restarts \emph{and} the best of a fixed number &
One algorithm ranks first under the former and last under the latter; two gains
near $+0.15$ become losses (\S\ref{sec:protocol:compute})\\
Degree mechanics read as structure &
Degree-preserving rewiring &
Both quantities are reproduced by the null, and are larger on it in 15 of 17 and
16 of 17 networks (\S\ref{sec:protocol:nulls})\\
Weighting credited to similarity &
Permute the weights across edges &
The permuted weights reproduce the gain and exceed it, $+0.0117$ against
$+0.0102$ (\S\ref{sec:protocol:nulls})\\
Nominal retention &
Report and match on measured retention &
The same nominal $\alpha$ realizes retentions from $0.33$ to $0.95$ depending on
the sampler (\S\ref{sec:protocol:retention})\\
Sparsification charged to nobody &
Time the pipeline, sparsifier included &
No quality-preserving configuration exceeds $0.66\times$
(\S\ref{sec:protocol:cost})\\
\bottomrule
\end{tabular}
\end{table*}

\subsection{Pre-registration}\label{sec:protocol:prereg}

Each experiment in this study has a design document written before its code,
fixing what is measured, what is predicted, and what outcome would end the line of
work. Several of the predictions were ours and failed. The synthetic sweep that
locates the boundary was registered with the prediction that sparsification would
not produce an honest modularity gain at any density, and the fixed-$k$ arm
falsified it.
